\documentclass[10pt]{article}
% Hojas resumen: un A4 por delante y por detrás, dos columnas por cara.
% Uso: \documentclass[10pt]{article} % Hojas resumen: un A4 por delante y por detrás, dos columnas por cara.
% Uso: \documentclass[10pt]{article} % Hojas resumen: un A4 por delante y por detrás, dos columnas por cara.
% Uso: \documentclass[10pt]{article} \input{estilo-resumen.tex} \bloque{analisis|algebra|azar} \cabecera{n}{Título}{Lecciones}
\usepackage[utf8]{inputenc}
\usepackage[T1]{fontenc}
\usepackage[spanish,es-nodecimaldot]{babel}
\usepackage{lmodern}
\usepackage{amsmath,amssymb}
\usepackage{graphicx}
\usepackage[a4paper,left=10mm,right=10mm,top=15mm,bottom=13mm,headheight=11pt,headsep=3mm,footskip=6mm]{geometry}
\usepackage{multicol}
\setlength{\columnsep}{6mm}
\setlength{\columnseprule}{0.3pt}
\usepackage{xcolor}
\usepackage{enumitem}
\setlist{nosep,leftmargin=1.1em,itemsep=1pt}
\usepackage{array}
\usepackage{needspace}
\usepackage{pgfplots}
\usepgfplotslibrary{fillbetween}
\pgfplotsset{compat=1.18,
  mini/.style={width=4.3cm,height=3.3cm,scale only axis=false,axis lines=middle,xtick=\empty,ytick=\empty,samples=120,
    xlabel={\scriptsize$x$},ylabel={\scriptsize$y$},title style={font=\scriptsize,yshift=-4pt},clip=true,
    every axis plot/.append style={thick,color=acento}}}
\usepackage{ragged2e}
\usepackage[most]{tcolorbox}
\usepackage{fancyhdr,lastpage}
\usepackage[hidelinks,pdfencoding=auto]{hyperref}
\hypersetup{pdflang=es-ES,pdfsubject={Matemáticas Aplicadas a las Ciencias Sociales II}}

\definecolor{analisis}{HTML}{1D4ED8}
\definecolor{algebra}{HTML}{86198F}
\definecolor{azar}{HTML}{9A3412}
\definecolor{aviso}{HTML}{B91C1C}
\colorlet{acento}{analisis}
\newcommand{\bloque}[1]{\colorlet{acento}{#1}}

\setlength{\parindent}{0pt}
\setlength{\parskip}{2pt}
\renewcommand{\arraystretch}{1.2}
\AtBeginDocument{\RaggedRight\fontsize{9.8}{11.9}\selectfont\setlength{\abovedisplayskip}{3pt}\setlength{\belowdisplayskip}{3pt}\setlength{\abovedisplayshortskip}{1pt}\setlength{\belowdisplayshortskip}{1pt}}

\newcommand{\temanum}{}\newcommand{\tematitulo}{}
\pagestyle{fancy}
\fancyhf{}
\fancyhead[L]{\footnotesize\color{acento}\textbf{Hoja resumen}\quad Tema \temanum. \tematitulo}
\fancyhead[R]{\footnotesize Matemáticas Aplicadas a las Ciencias Sociales II --- 2.\textsuperscript{o} Bachillerato}
\fancyfoot[L]{\scriptsize IES Al-Ándalus --- Juan María Gámez Ortiz, con la ayuda de Claude y de Codex --- CC BY-NC 4.0}
\fancyfoot[R]{\scriptsize Cara \thepage\ de \pageref{LastPage}}
\renewcommand{\headrulewidth}{0.4pt}

% Cabecera de la primera cara: número, título y lecciones del tema.
\newcommand{\cabecera}[3]{%
  \renewcommand{\temanum}{#1}\renewcommand{\tematitulo}{#2}%
  \hypersetup{pdftitle={Hoja resumen. Tema #1. #2}}%
  \noindent\begin{tcolorbox}[enhanced,colback=acento!8,colframe=acento,boxrule=0.6pt,arc=2pt,left=4pt,right=4pt,top=2pt,bottom=2pt]
  \begin{minipage}[c]{1.1cm}\includegraphics[width=1cm]{logo-al-andalus.png}\end{minipage}%
  \begin{minipage}[c]{\dimexpr\linewidth-1.2cm}{\large\bfseries\color{acento}Tema #1. #2}\\[1pt]{\scriptsize #3}\end{minipage}
  \end{tcolorbox}\vspace{-2pt}}

% Apartado con filete del color del bloque.
\newcommand{\seccion}[1]{\par\needspace{5\baselineskip}\vspace{3pt}{\color{acento}\bfseries\normalsize #1}\par\vspace{-1pt}{\color{acento}\hrule height 0.5pt}\vspace{2pt}}
\newcommand{\sub}[1]{\par\needspace{3\baselineskip}\vspace{1pt}\textbf{\color{acento}#1}\enspace}

% Fórmula clave, aviso de errores y ejemplo.
\tcbset{before upper=\RaggedRight}
\newtcolorbox{formula}{enhanced,colback=acento!6,colframe=acento!70,boxrule=0.5pt,arc=1.5pt,left=3pt,right=3pt,top=0pt,bottom=0pt,before skip=2pt,after skip=2pt}
\newtcolorbox{ojo}{enhanced,colback=aviso!5,colframe=aviso!60,boxrule=0.5pt,arc=1.5pt,left=3pt,right=3pt,top=1pt,bottom=1pt,before skip=3pt,after skip=2pt,
  fonttitle=\bfseries\footnotesize,coltitle=aviso,colbacktitle=aviso!5,title={Ojo},attach title to upper={\enspace}}
\newtcolorbox{ejemplo}{enhanced,colback=white,colframe=acento!45,boxrule=0.4pt,arc=1.5pt,left=3pt,right=3pt,top=1pt,bottom=1pt,before skip=3pt,after skip=2pt,
  fonttitle=\bfseries\footnotesize,coltitle=acento,colbacktitle=white,title={Ejemplo},attach title to upper={\enspace}}
\newcommand{\dx}{\,dx}
 \bloque{analisis|algebra|azar} \cabecera{n}{Título}{Lecciones}
\usepackage[utf8]{inputenc}
\usepackage[T1]{fontenc}
\usepackage[spanish,es-nodecimaldot]{babel}
\usepackage{lmodern}
\usepackage{amsmath,amssymb}
\usepackage{graphicx}
\usepackage[a4paper,left=10mm,right=10mm,top=15mm,bottom=13mm,headheight=11pt,headsep=3mm,footskip=6mm]{geometry}
\usepackage{multicol}
\setlength{\columnsep}{6mm}
\setlength{\columnseprule}{0.3pt}
\usepackage{xcolor}
\usepackage{enumitem}
\setlist{nosep,leftmargin=1.1em,itemsep=1pt}
\usepackage{array}
\usepackage{needspace}
\usepackage{pgfplots}
\usepgfplotslibrary{fillbetween}
\pgfplotsset{compat=1.18,
  mini/.style={width=4.3cm,height=3.3cm,scale only axis=false,axis lines=middle,xtick=\empty,ytick=\empty,samples=120,
    xlabel={\scriptsize$x$},ylabel={\scriptsize$y$},title style={font=\scriptsize,yshift=-4pt},clip=true,
    every axis plot/.append style={thick,color=acento}}}
\usepackage{ragged2e}
\usepackage[most]{tcolorbox}
\usepackage{fancyhdr,lastpage}
\usepackage[hidelinks,pdfencoding=auto]{hyperref}
\hypersetup{pdflang=es-ES,pdfsubject={Matemáticas Aplicadas a las Ciencias Sociales II}}

\definecolor{analisis}{HTML}{1D4ED8}
\definecolor{algebra}{HTML}{86198F}
\definecolor{azar}{HTML}{9A3412}
\definecolor{aviso}{HTML}{B91C1C}
\colorlet{acento}{analisis}
\newcommand{\bloque}[1]{\colorlet{acento}{#1}}

\setlength{\parindent}{0pt}
\setlength{\parskip}{2pt}
\renewcommand{\arraystretch}{1.2}
\AtBeginDocument{\RaggedRight\fontsize{9.8}{11.9}\selectfont\setlength{\abovedisplayskip}{3pt}\setlength{\belowdisplayskip}{3pt}\setlength{\abovedisplayshortskip}{1pt}\setlength{\belowdisplayshortskip}{1pt}}

\newcommand{\temanum}{}\newcommand{\tematitulo}{}
\pagestyle{fancy}
\fancyhf{}
\fancyhead[L]{\footnotesize\color{acento}\textbf{Hoja resumen}\quad Tema \temanum. \tematitulo}
\fancyhead[R]{\footnotesize Matemáticas Aplicadas a las Ciencias Sociales II --- 2.\textsuperscript{o} Bachillerato}
\fancyfoot[L]{\scriptsize IES Al-Ándalus --- Juan María Gámez Ortiz, con la ayuda de Claude y de Codex --- CC BY-NC 4.0}
\fancyfoot[R]{\scriptsize Cara \thepage\ de \pageref{LastPage}}
\renewcommand{\headrulewidth}{0.4pt}

% Cabecera de la primera cara: número, título y lecciones del tema.
\newcommand{\cabecera}[3]{%
  \renewcommand{\temanum}{#1}\renewcommand{\tematitulo}{#2}%
  \hypersetup{pdftitle={Hoja resumen. Tema #1. #2}}%
  \noindent\begin{tcolorbox}[enhanced,colback=acento!8,colframe=acento,boxrule=0.6pt,arc=2pt,left=4pt,right=4pt,top=2pt,bottom=2pt]
  \begin{minipage}[c]{1.1cm}\includegraphics[width=1cm]{logo-al-andalus.png}\end{minipage}%
  \begin{minipage}[c]{\dimexpr\linewidth-1.2cm}{\large\bfseries\color{acento}Tema #1. #2}\\[1pt]{\scriptsize #3}\end{minipage}
  \end{tcolorbox}\vspace{-2pt}}

% Apartado con filete del color del bloque.
\newcommand{\seccion}[1]{\par\needspace{5\baselineskip}\vspace{3pt}{\color{acento}\bfseries\normalsize #1}\par\vspace{-1pt}{\color{acento}\hrule height 0.5pt}\vspace{2pt}}
\newcommand{\sub}[1]{\par\needspace{3\baselineskip}\vspace{1pt}\textbf{\color{acento}#1}\enspace}

% Fórmula clave, aviso de errores y ejemplo.
\tcbset{before upper=\RaggedRight}
\newtcolorbox{formula}{enhanced,colback=acento!6,colframe=acento!70,boxrule=0.5pt,arc=1.5pt,left=3pt,right=3pt,top=0pt,bottom=0pt,before skip=2pt,after skip=2pt}
\newtcolorbox{ojo}{enhanced,colback=aviso!5,colframe=aviso!60,boxrule=0.5pt,arc=1.5pt,left=3pt,right=3pt,top=1pt,bottom=1pt,before skip=3pt,after skip=2pt,
  fonttitle=\bfseries\footnotesize,coltitle=aviso,colbacktitle=aviso!5,title={Ojo},attach title to upper={\enspace}}
\newtcolorbox{ejemplo}{enhanced,colback=white,colframe=acento!45,boxrule=0.4pt,arc=1.5pt,left=3pt,right=3pt,top=1pt,bottom=1pt,before skip=3pt,after skip=2pt,
  fonttitle=\bfseries\footnotesize,coltitle=acento,colbacktitle=white,title={Ejemplo},attach title to upper={\enspace}}
\newcommand{\dx}{\,dx}
 \bloque{analisis|algebra|azar} \cabecera{n}{Título}{Lecciones}
\usepackage[utf8]{inputenc}
\usepackage[T1]{fontenc}
\usepackage[spanish,es-nodecimaldot]{babel}
\usepackage{lmodern}
\usepackage{amsmath,amssymb}
\usepackage{graphicx}
\usepackage[a4paper,left=10mm,right=10mm,top=15mm,bottom=13mm,headheight=11pt,headsep=3mm,footskip=6mm]{geometry}
\usepackage{multicol}
\setlength{\columnsep}{6mm}
\setlength{\columnseprule}{0.3pt}
\usepackage{xcolor}
\usepackage{enumitem}
\setlist{nosep,leftmargin=1.1em,itemsep=1pt}
\usepackage{array}
\usepackage{needspace}
\usepackage{pgfplots}
\usepgfplotslibrary{fillbetween}
\pgfplotsset{compat=1.18,
  mini/.style={width=4.3cm,height=3.3cm,scale only axis=false,axis lines=middle,xtick=\empty,ytick=\empty,samples=120,
    xlabel={\scriptsize$x$},ylabel={\scriptsize$y$},title style={font=\scriptsize,yshift=-4pt},clip=true,
    every axis plot/.append style={thick,color=acento}}}
\usepackage{ragged2e}
\usepackage[most]{tcolorbox}
\usepackage{fancyhdr,lastpage}
\usepackage[hidelinks,pdfencoding=auto]{hyperref}
\hypersetup{pdflang=es-ES,pdfsubject={Matemáticas Aplicadas a las Ciencias Sociales II}}

\definecolor{analisis}{HTML}{1D4ED8}
\definecolor{algebra}{HTML}{86198F}
\definecolor{azar}{HTML}{9A3412}
\definecolor{aviso}{HTML}{B91C1C}
\colorlet{acento}{analisis}
\newcommand{\bloque}[1]{\colorlet{acento}{#1}}

\setlength{\parindent}{0pt}
\setlength{\parskip}{2pt}
\renewcommand{\arraystretch}{1.2}
\AtBeginDocument{\RaggedRight\fontsize{9.8}{11.9}\selectfont\setlength{\abovedisplayskip}{3pt}\setlength{\belowdisplayskip}{3pt}\setlength{\abovedisplayshortskip}{1pt}\setlength{\belowdisplayshortskip}{1pt}}

\newcommand{\temanum}{}\newcommand{\tematitulo}{}
\pagestyle{fancy}
\fancyhf{}
\fancyhead[L]{\footnotesize\color{acento}\textbf{Hoja resumen}\quad Tema \temanum. \tematitulo}
\fancyhead[R]{\footnotesize Matemáticas Aplicadas a las Ciencias Sociales II --- 2.\textsuperscript{o} Bachillerato}
\fancyfoot[L]{\scriptsize IES Al-Ándalus --- Juan María Gámez Ortiz, con la ayuda de Claude y de Codex --- CC BY-NC 4.0}
\fancyfoot[R]{\scriptsize Cara \thepage\ de \pageref{LastPage}}
\renewcommand{\headrulewidth}{0.4pt}

% Cabecera de la primera cara: número, título y lecciones del tema.
\newcommand{\cabecera}[3]{%
  \renewcommand{\temanum}{#1}\renewcommand{\tematitulo}{#2}%
  \hypersetup{pdftitle={Hoja resumen. Tema #1. #2}}%
  \noindent\begin{tcolorbox}[enhanced,colback=acento!8,colframe=acento,boxrule=0.6pt,arc=2pt,left=4pt,right=4pt,top=2pt,bottom=2pt]
  \begin{minipage}[c]{1.1cm}\includegraphics[width=1cm]{logo-al-andalus.png}\end{minipage}%
  \begin{minipage}[c]{\dimexpr\linewidth-1.2cm}{\large\bfseries\color{acento}Tema #1. #2}\\[1pt]{\scriptsize #3}\end{minipage}
  \end{tcolorbox}\vspace{-2pt}}

% Apartado con filete del color del bloque.
\newcommand{\seccion}[1]{\par\needspace{5\baselineskip}\vspace{3pt}{\color{acento}\bfseries\normalsize #1}\par\vspace{-1pt}{\color{acento}\hrule height 0.5pt}\vspace{2pt}}
\newcommand{\sub}[1]{\par\needspace{3\baselineskip}\vspace{1pt}\textbf{\color{acento}#1}\enspace}

% Fórmula clave, aviso de errores y ejemplo.
\tcbset{before upper=\RaggedRight}
\newtcolorbox{formula}{enhanced,colback=acento!6,colframe=acento!70,boxrule=0.5pt,arc=1.5pt,left=3pt,right=3pt,top=0pt,bottom=0pt,before skip=2pt,after skip=2pt}
\newtcolorbox{ojo}{enhanced,colback=aviso!5,colframe=aviso!60,boxrule=0.5pt,arc=1.5pt,left=3pt,right=3pt,top=1pt,bottom=1pt,before skip=3pt,after skip=2pt,
  fonttitle=\bfseries\footnotesize,coltitle=aviso,colbacktitle=aviso!5,title={Ojo},attach title to upper={\enspace}}
\newtcolorbox{ejemplo}{enhanced,colback=white,colframe=acento!45,boxrule=0.4pt,arc=1.5pt,left=3pt,right=3pt,top=1pt,bottom=1pt,before skip=3pt,after skip=2pt,
  fonttitle=\bfseries\footnotesize,coltitle=acento,colbacktitle=white,title={Ejemplo},attach title to upper={\enspace}}
\newcommand{\dx}{\,dx}

\bloque{algebra}
\begin{document}
\cabecera{4}{Matrices y determinantes}{Lecciones 4.1--4.12: tipos de matrices, operaciones, producto, potencias, traspuesta y simetría, determinantes y sus propiedades, desarrollo por adjuntos, matriz inversa, rango, ecuaciones matriciales y grafos.}
\begin{multicols}{2}

\seccion{1. Matrices: definición y tipos}
Una matriz \(A=(a_{ij})\) de dimensión \(m\times n\) tiene \(m\) filas y \(n\) columnas; \(a_{ij}\) está en la fila \(i\) y la columna \(j\).
\begin{itemize}
  \item \textbf{Fila} \(1\times n\), \textbf{columna} \(m\times1\), \textbf{cuadrada} \(n\times n\) (diagonal principal \(a_{11},\dots,a_{nn}\)).
  \item \textbf{Nula} \(O\); \textbf{diagonal} (ceros fuera de la diagonal); \textbf{identidad} \(I\) (diagonal de unos).
  \item \textbf{Triangular} superior o inferior (ceros debajo o encima de la diagonal).
  \item \textbf{Traspuesta} \(A^T\): las filas pasan a columnas, \((A^T)_{ij}=a_{ji}\).
  \item \textbf{Simétrica} \(A^T=A\). \textbf{Antisimétrica} \(A^T=-A\) (diagonal nula).
\end{itemize}
Dos matrices son iguales si tienen la misma dimensión y los mismos elementos en cada posición.

\seccion{2. Operaciones}
\sub{Suma y producto por un número:} elemento a elemento; la suma exige la misma dimensión.
\begin{formula}
\[\underbrace{A}_{m\times n}\ \underbrace{B}_{n\times p}=\underbrace{C}_{m\times p}\]
\[c_{ij}=\sum_{k=1}^{n}a_{ik}b_{kj}\]
\end{formula}
\(c_{ij}\) = fila \(i\) de \(A\) por columna \(j\) de \(B\). Solo existe si las columnas de \(A\) coinciden con las filas de \(B\).
\begin{ejemplo}
\[AB=\begin{pmatrix}1&2\\0&1\end{pmatrix}\begin{pmatrix}3&1\\1&0\end{pmatrix}=\begin{pmatrix}5&1\\1&0\end{pmatrix}\]
\[BA=\begin{pmatrix}3&1\\1&0\end{pmatrix}\begin{pmatrix}1&2\\0&1\end{pmatrix}=\begin{pmatrix}3&7\\1&2\end{pmatrix}\]
\end{ejemplo}
\begin{ojo}
\(AB\ne BA\) en general. Por eso \((A+B)^2\ne A^2+2AB+B^2\) y \(AB=O\) no implica \(A=O\) o \(B=O\).
\end{ojo}
\sub{Traspuesta.}
\[(A^T)^T=A,\qquad (A+B)^T=A^T+B^T\]
\[(\alpha A)^T=\alpha A^T,\qquad (AB)^T=B^TA^T\]

\seccion{3. Potencias}
\(A^n=A\cdot A\cdots A\) (solo matrices cuadradas). Se calculan \(A^2\), \(A^3\)\ldots{} hasta ver un \textbf{patrón} o una repetición, y se comprueba.
\begin{ejemplo}
\(A=\begin{pmatrix}1&1\\0&1\end{pmatrix}\Rightarrow A^2=\begin{pmatrix}1&2\\0&1\end{pmatrix},\ A^n=\begin{pmatrix}1&n\\0&1\end{pmatrix}\).
Si \(A^2=I\), entonces \(A^{n}=I\) para \(n\) par y \(A^n=A\) para \(n\) impar.
\end{ejemplo}

\seccion{4. Determinantes}
\begin{formula}
\[\begin{vmatrix}a&b\\c&d\end{vmatrix}=ad-bc\]
\[\begin{vmatrix}a&b&c\\d&e&f\\g&h&i\end{vmatrix}=aei+bfg+cdh\]
\[\hphantom{\begin{vmatrix}a\\d\\g\end{vmatrix}}\qquad\quad-(ceg+afh+bdi)\]
\end{formula}
\textbf{Regla de Sarrus:} tres productos con signo \(+\) (diagonal principal y sus paralelas) y tres con signo \(-\) (diagonal secundaria y sus paralelas).
\begin{ejemplo}
\[\begin{vmatrix}1&2&0\\0&1&1\\1&0&1\end{vmatrix}=(1+2+0)-(0+0+0)=3\]
\end{ejemplo}

\seccion{5. Propiedades de los determinantes}
\begin{formula}
\[\det(A^T)=\det(A),\qquad \det(AB)=\det(A)\det(B)\]
\[\det(kA)=k^n\det(A)\quad(A\text{ de orden }n)\]
\end{formula}
\begin{itemize}
  \item Una fila (o columna) de ceros, o dos filas iguales o proporcionales: \(\det=0\).
  \item Intercambiar dos filas cambia el signo.
  \item Multiplicar una fila por \(k\) multiplica el determinante por \(k\).
  \item Sumar a una fila una combinación de otras no lo cambia (así se hacen ceros).
  \item Matriz triangular: producto de la diagonal.
\end{itemize}

\seccion{6. Desarrollo por adjuntos}
\begin{formula}
\[M_{ij}=\text{menor complementario},\qquad C_{ij}=(-1)^{i+j}M_{ij}\]
\[\det(A)=\sum_{j=1}^{n}a_{ij}C_{ij}\quad\text{(por la fila }i\text{)}\]
\end{formula}
\(M_{ij}\) es el determinante que queda al suprimir la fila \(i\) y la columna \(j\). Signos \(\begin{smallmatrix}+&-&+\\-&+&-\\+&-&+\end{smallmatrix}\). Conviene desarrollar por la fila o columna con más ceros (o crearlos antes).

\seccion{7. Matriz inversa}
\(A\) es \textbf{invertible} (regular) \(\iff\det(A)\ne0\). Entonces \(AA^{-1}=A^{-1}A=I\).
\begin{formula}
\[A^{-1}=\frac1{\det(A)}\operatorname{Adj}(A),\qquad \operatorname{Adj}(A)=C^T\]
\[A=\begin{pmatrix}a&b\\c&d\end{pmatrix}\ \Rightarrow\]
\[A^{-1}=\frac1{ad-bc}\begin{pmatrix}d&-b\\-c&a\end{pmatrix}\]
\end{formula}
\textbf{Pasos:} 1) \(\det(A)\ne0\); 2) matriz de adjuntos \(C\); 3) trasponer: \(\operatorname{Adj}(A)=C^T\); 4) dividir entre \(\det(A)\); 5) comprobar \(AA^{-1}=I\).
\begin{ejemplo}
\(A=\begin{pmatrix}2&1\\1&2\end{pmatrix}\): \(\det A=3\), \(A^{-1}=\dfrac13\begin{pmatrix}2&-1\\-1&2\end{pmatrix}\).
\end{ejemplo}
Propiedades: \((AB)^{-1}=B^{-1}A^{-1}\), \((A^{-1})^{-1}=A\), \(\det(A^{-1})=\dfrac1{\det A}\).

\seccion{8. Rango}
\begin{formula}
\[0\le\operatorname{rg}(A)\le\min(m,n)\]
\end{formula}
El rango es el número máximo de filas (o columnas) linealmente independientes, o el orden del mayor menor no nulo. También: número de filas no nulas tras escalonar por Gauss.
\begin{ejemplo}
\(A=\begin{pmatrix}1&2&3\\2&4&6\\1&0&1\end{pmatrix}\): la fila 2 es el doble de la 1, así que \(\det A=0\); el menor \(\begin{vmatrix}1&2\\1&0\end{vmatrix}=-2\ne0\Rightarrow\operatorname{rg}(A)=2\).
\end{ejemplo}

\seccion{9. Ecuaciones matriciales}
\begin{formula}
\[AX=B\Rightarrow X=A^{-1}B,\qquad XA=B\Rightarrow X=BA^{-1}\]
\end{formula}
Se despeja como con números, pero \textbf{respetando el lado}: se multiplica por \(A^{-1}\) por la izquierda o por la derecha. \(AX+B=C\Rightarrow X=A^{-1}(C-B)\).
\begin{ejemplo}
\(A=\begin{pmatrix}2&1\\1&1\end{pmatrix}\), \(B=\begin{pmatrix}3\\2\end{pmatrix}\): \(\det A=1\), \(A^{-1}=\begin{pmatrix}1&-1\\-1&2\end{pmatrix}\), \(X=A^{-1}B=\begin{pmatrix}1\\1\end{pmatrix}\). Comprobación: \(AX=B\).
\end{ejemplo}

\seccion{10. Grafos y matriz de adyacencia}
\(a_{ij}=1\) si hay arista (o arco) del vértice \(i\) al \(j\), y 0 si no. Grafo no dirigido: matriz simétrica.
\begin{formula}
\((A^k)_{ij}\) = número de recorridos de exactamente \(k\) pasos del vértice \(i\) al \(j\).
\end{formula}
\begin{ejemplo}
Vértices 1, 2, 3 con aristas 1--2 y 2--3: \(A=\begin{pmatrix}0&1&0\\1&0&1\\0&1&0\end{pmatrix}\), \(A^2=\begin{pmatrix}1&0&1\\0&2&0\\1&0&1\end{pmatrix}\). Hay 2 recorridos de dos pasos de 2 a 2 (por 1 y por 3).
\end{ejemplo}

\seccion{11. Inversa de orden 3 paso a paso}
\(A=\begin{pmatrix}1&2&0\\0&1&1\\1&0&1\end{pmatrix}\), \(\det A=3\ne0\).
\begin{enumerate}
  \item Adjuntos: \(C_{11}=1,\ C_{12}=1,\ C_{13}=-1\); \(C_{21}=-2,\ C_{22}=1,\ C_{23}=2\); \(C_{31}=2,\ C_{32}=-1,\ C_{33}=1\).
  \item Trasponer y dividir entre 3:
\end{enumerate}
\[A^{-1}=\frac13\begin{pmatrix}1&-2&2\\1&1&-1\\-1&2&1\end{pmatrix}\]
Comprobación: \(A\cdot\operatorname{Adj}(A)=3I\), así que \(AA^{-1}=I\).

\seccion{12. Matrices con un parámetro}
Se estudia \(\det\) en función del parámetro: los valores que lo anulan son los casos especiales.
\begin{ejemplo}
\(A=\begin{pmatrix}1&1\\1&k\end{pmatrix}\): \(\det A=k-1\).
Si \(k\ne1\): \(\operatorname{rg}(A)=2\) y \(A\) es invertible.
Si \(k=1\): filas iguales, \(\operatorname{rg}(A)=1\) y no tiene inversa.
\end{ejemplo}

\seccion{13. Matrices en contexto}
Una matriz ordena datos: filas y columnas con significado. El producto combina tablas cuando las unidades encajan.
\begin{ejemplo}
Dos tiendas venden tres productos. Unidades \(Q=\begin{pmatrix}10&5&2\\4&8&6\end{pmatrix}\) (filas: tiendas) y precios \(P=\begin{pmatrix}2\\3\\5\end{pmatrix}\) euros (por producto). \(QP=\begin{pmatrix}45\\62\end{pmatrix}\): ingresos de cada tienda, en euros.
\end{ejemplo}

\begin{ojo}
\begin{itemize}
  \item Multiplicar matrices elemento a elemento.
  \item Escribir \(X=BA^{-1}\) cuando la ecuación es \(AX=B\).
  \item Olvidar trasponer la matriz de adjuntos o el signo \((-1)^{i+j}\).
  \item \(\det(kA)=k\det(A)\): falta la potencia \(k^n\).
  \item Calcular \(A^{-1}\) sin comprobar que \(\det A\ne0\).
\end{itemize}
\end{ojo}

\end{multicols}
\end{document}
